\chapter{Theory}
\minitoc

\section{Trapping Ions}
\label{sec:Intro Ion trap}

\subsection{Trapping Potential}
% Main points:
    % Give useful equations for describing trapping pseudopotential
    % Give ion spacing and mode frequencies
    % Could put tickle interaction here
% Pre reqs:
    % 
    From Earnshaw's theorem,
    $\nabla^2 V = 0,$ a stable stationary point in 3D cannot be realised
    using only static electric potentials, $V$, as if the potential is confining
    in two dimensions, it will be anti-confining in the third. Therefore, to
    achieve stable trapping, either an oscillating electric field (Paul
    trap~\cite{paul1990electromagnetic}), or a static magnetic field (Penning trap) is
    used.\\
Will include pseudopotential relating $V_{RF}$ and $\Omega_{RF}$ to the radial mode frequencies $\omega_m$ (see Vera eqn 3.4),
\begin{equation}\label{eqn:pseudo_pot}
    \omega_r \propto \frac{V_{RF}}{M\Omega_{RF}}.
\end{equation}
Further, will use James to give ion spacing and non-COM mode frequencies as will be relevant for single addressing, fast gates (if discussed), and for motional mode excitati as will be relevant for single addressing, fast gates (if discussed), and for motional mode excitation.

\section{Light Matter Interactions}
% Main points:
    % Give useful equations for describing Rabi flopping
    % For motional mode coupling
    % SDF
% Pre reqs:
    % 

\subsection{Rabi Flopping}
\label{sec:Rabi Flopping}
\subsection{Motional Coupling}
\label{sec:Motional Coupling}
% Sidebands
% SDFs
\subsection{Geometric Phase Gates}
\label{sec:Geometric Phase Gates}
%
% Overlapped theory for chapters \ref{ch:virtual_distillation} and \ref{ch:rb} 
\section{Randomised Benchmarking}
Some overview and motivation for RB.
    Need to cover:
    \begin{itemize}
        \item Twirling over a group
        \item Liouville representation of a CPTP map
        \item Irreps of the Clifford group
        \item Clifford approximates twirling over the unitary group (Haar measure)
        \item How Liouville representation is transformed under twirling
        \item The specific representation-theory tools you use later: irreps, multiplicity spaces, invariant sectors, Schur basis.\\
        \item RB protocol for single qubit Clifford group, generalisation to multi-qubit Clifford group.
        \item How to extract error per Clifford from decay rate
\end{itemize}

\section{Noise channels}
    As will also be used in virtual distillation, may move to theory sections\\
    A description of the noise channels we consider in this work. Explicit CPTP maps for these channels.
    \subsection{Why isn't everything white noise?}
    \subsection{Unitary noise}
        Both continuous and probabilistic.
    \subsection{Depolarising noise}
        An ensemble of unitary noise channels, with a uniform distribution over the Bloch sphere.
        Get this from twirling over Clifford group!
    \subsection{Spatially correlated noise}
        \indent\textbf{Common noise}
            An ensemble of unitary noise channels, with a uniform distribution over the Bloch sphere, but the same unitary is applied to multiple qubits. Can enforce some symmetry to this channel by twirling over the k-copy Clifford group. This is the channel we will use to model correlated noise in our experiments.
        \indent\textbf{Entangling noise}
            In SC ZZ entangling noise is common.
    \subsection{Can we make everything white noise?}
        Have a look at Balint paper, deep random circuits, etc.

\section{v1}

This chapter lays the theoretical groundwork for Chapter~\ref{chap:results}.
It reviews standard Clifford randomised benchmarking (RB), develops the
representation-theoretic machinery of group twirling in some generality
before specialising it to the two-copy Clifford group, and surveys the
classes of noise channel and prior characterisation strategies that motivate
a symmetry-based approach to correlated-noise detection.

% ==========================================================================
\section{Standard Clifford randomised benchmarking}
\label{sec:standard_rb}

\subsection{Motivation: characterising noise without full tomography}

As quantum processors scale, fault tolerance and error mitigation become
necessary to overcome decoherence and systematic control
errors~\cite{shor_fault-tolerant_1996, aharonov_fault-tolerant_2008,
cai_quantum_2023}. Errors that induce correlations across multiple qubits ---
variously termed spatial correlations or crosstalk
errors~\cite{sarovar_detecting_2020} --- require that mitigation and
correction strategies be built on noise models that faithfully capture their
correlated structure. Assuming a globally depolarising (white) noise model
where correlations are in fact present can lead to systematic
underestimation of logical error rates~\cite{aharonov_fault-tolerant_2006},
or systematic bias in the outcome of an error-corrected
observable~\cite{foldager_can_2024}.

Methods exist to simplify correlated noise channels into more favourable
depolarising or white-noise forms, for example Pauli twirling and random
circuit compilation~\cite{geller_efficient_2013, wallman_noise_2016,
hashim_benchmarking_2023}. However, the efficacy of such methods must still
be characterised, verified, and optimised on physical devices with realistic
noise, which is only possible with a scalable diagnostic for the correlated
component of that noise in the first place.

\todo{Add a short paragraph situating this thesis chapter within the broader
RB literature landscape (interleaved RB, cycle benchmarking, purity
benchmarking, etc.), with a comparison table of scaling/expressiveness
trade-offs. This context was compressed to a single sentence in the PRL and
would benefit from a fuller literature review here.}

\subsection{The single-qubit Clifford RB protocol}

Standard single-qubit Clifford RB proceeds by applying a sequence of $M-1$
Cliffords, sampled uniformly at random from the single-qubit Clifford group,
followed by a single inverting Clifford, $C_{\mathrm{inv}}$, chosen so that
the sequence implements the identity operation in the absence of
errors~\cite{dankert_exact_2009, magesan_scalable_2011}. For each sequence
length $M$, $r$ independent random sequences are generated, and each is
repeated $s$ times to estimate the probability of the qubit returning to its
target computational state. The target state is conventionally randomised
between $\ket{0}$ and $\ket{1}$ from sequence to sequence, which reduces
sensitivity to asymmetries in state preparation and
measurement~\cite{knill_randomized_2008}.

\todo{Circuit diagram: single-qubit Clifford RB sequence of length
$M$, with per-step error box and final inverting Clifford $C_{\mathrm{inv}}$.}
{Schematic single-qubit Clifford RB circuit (reproduced/adapted from
Fig.~1(a) of the manuscript). \todo{Insert final figure once redrawn for
thesis formatting; the manuscript version also shows $k$ parallel copies,
which belongs in Chapter~\ref{chap:results} instead.}}

The survival probability of a single qubit after $m$ Cliffords follows a
single exponential decay,
\begin{equation}
    \mathcal{P}_0 = \frac{1}{2}\left(1 + (1-\varepsilon_c)^m\right),
    \label{eqn:1q_rb_decay}
\end{equation}
where $\varepsilon_c$ is the ``error per Clifford'' (EPC), the single
scalar figure of merit that standard RB is able to extract. This compression
of an entire error channel into one averaged parameter is precisely what
gives RB its scalability, but it is also the limitation that motivates the
correlated-noise-resolving protocol developed in
Chapter~\ref{chap:results}: two qubits with identical single-qubit EPC can
have entirely different amounts of correlated error between them, and
Eq.~\eqref{eqn:1q_rb_decay} cannot distinguish these cases.

\subsection{State-preparation and measurement (SPAM) errors}
\label{sec:background_spam}

RB's scalability partly derives from its insensitivity to SPAM errors up to
an overall scaling of the exponential decay. Modelling independent bit-flip
type state-preparation error on a single qubit with probability
$\varepsilon_{sp}^{(i)}/2$, the prepared state is
\begin{equation}
    \rho_{sp}^{(i)} = \left(1-\frac{\varepsilon_{sp}^{(i)}}{2}\right)\ket{0}\bra{0}
    + \frac{\varepsilon_{sp}^{(i)}}{2}\ket{1}\bra{1},
\end{equation}
equivalently expressed as a probabilistic $X$ channel acting on an arbitrary
input state,
\begin{equation}
    \rho \mapsto \mathcal{E}_{sp}^{(i)}(\rho)
    = \left(1-\frac{\varepsilon_{sp}^{(i)}}{2}\right)\rho
    + \frac{\varepsilon_{sp}^{(i)}}{2} X\rho X.
    \label{eqn:spam_channel}
\end{equation}
In the Liouville picture this simply rescales the overlap of the initial
state with each sub-representation of the twirled channel by a factor
$A_{sp}^{(i)} = (1-\varepsilon_{sp}^{(i)})$ for every sub-representation in
which qubit $i$ participates non-trivially. This single-qubit treatment
generalises directly to the two-qubit setting used in
Chapter~\ref{chap:results}: correlated SPAM errors, such as a correlated
spin-flip channel, leave certain sub-representations (in particular the
fully symmetric sector) invariant, and so can mimic a genuine correlated
error signature if left unaccounted for in a fit.

\todo{Expand this subsection with the general $n$-qubit SPAM formalism, and
a short discussion of readout-error mitigation strategies compared against
this multiplicative-scaling picture. The manuscript SM only develops the
2-qubit special case, which is deferred to Chapter~\ref{chap:results}.}

% ==========================================================================
\section{Twirling and representation theory}
\label{sec:twirling_rep_theory}

\subsection{General twirling formalism}

RB owes its practicality to the process of \emph{twirling}: averaging a
general noise channel over a chosen group of unitaries. Consider a general
CPTP error channel $\Lambda$, expressed in Liouville (transfer-matrix) form
acting on a vectorised density matrix. Twirling $\Lambda$ over a discrete
group $\mathcal{G}$ produces
\begin{equation}
    \Lambda \mapsto \tilde\Lambda
    = \frac{1}{|\mathcal{G}|}\sum_{\mathcal{G}} G\Lambda G^\dagger,
    \label{eqn:twirl_general}
\end{equation}
where $G$ denotes the vectorised representation of a group element. By
construction, the twirled channel commutes with every representation
element,
\begin{equation}
    G\tilde\Lambda G^\dagger = \tilde\Lambda \qquad \forall\, G \in \mathcal{G}.
    \label{eqn:twirl_commute}
\end{equation}

The Hilbert space on which $\mathcal{G}$ acts decomposes into a direct sum
of irreducible representations (irreps),
\begin{equation}
    \mathcal{H} = \bigoplus_\lambda \left(V_\lambda \otimes \mathbb{C}^{m_\lambda}\right),
\end{equation}
where $V_\lambda$ is the irrep space and $\mathbb{C}^{m_\lambda}$ is a
multiplicity space of dimension $m_\lambda$. Schur's lemma then constrains
any operator satisfying Eq.~\eqref{eqn:twirl_commute} to be block diagonal
in this decomposition,
\begin{equation}
    \tilde\Lambda = \bigoplus_\lambda \left(I_\lambda \otimes \mathcal{A}_\lambda\right),
    \label{eqn:block_diagonal}
\end{equation}
where $I_\lambda$ is the identity on $V_\lambda$ and $\mathcal{A}_\lambda$ is
an $m_\lambda \times m_\lambda$ matrix acting only within the multiplicity
space. Physically, Eq.~\eqref{eqn:block_diagonal} is the statement that
twirling collapses an otherwise arbitrarily complex error channel down to a
comparatively small number of scalar decay rates (diagonal terms of
$\mathcal{A}_\lambda$) and couplings (off-diagonal terms), one set per irrep.
This is the general mechanism underlying the scalability of every RB variant,
and the choice of twirling group $\mathcal{G}$ determines which physical
information survives the averaging and which is discarded.

\todo{Add a short worked example here with a simple, well-known group (e.g.\
the single-qubit Clifford group acting on one qubit, recovering standard RB's
single decay parameter) to build intuition before jumping to the two-copy
case below. This would also be a natural place to discuss the trade-off
between expressiveness and scalability more formally, referencing
Pauli-twirling as an alternative choice of $\mathcal{G}$.}

\subsection{The two-copy Clifford group and its irreducible representations}
\label{sec:two_copy_cliff}

To characterise correlated noise between a pair of qubits, we are interested
in the group generated by applying the \emph{same} randomly sampled
single-qubit Clifford simultaneously to both qubits: the diagonal action of
the Clifford group on two copies, denoted $Cl^{\otimes 2}$. The
representation theory of this group has been analysed in detail
elsewhere~\cite{helsen_representations_2018, wallman_randomized_2014}; the
key results are summarised here.

The general $16$-dimensional two-qubit Liouville space decomposes under
$Cl^{\otimes 2}$ into four irreducible representations, of dimensions
$\{1, 2, 3, 3'\}$ (the label $3'$ distinguishes the second, inequivalent,
three-dimensional irrep from the first), with respective multiplicities
$\{2, 1, 3, 1\}$. Table~\ref{table:irreps} lists the six sub-representations
that follow from these irreps and multiplicities, together with a spanning
set of Pauli-basis operators for each.

\begin{table}[t]
    \centering
    \begin{tabular}{c|c|c}
         Sub-representation & irrep & Basis operators\\
         \hline\hline
         Identity & $1$ & $II$ \\
         \hline
         Left qubit & $3$ & $XI,YI,ZI$ \\
         \hline
         Right qubit & $3$ & $IX,IY,IZ$ \\
         \hline
         Trace & $1$ & $(XX + YY + ZZ)/\sqrt{3}$ \\
         \hline
         Diagonal traceless & $2$ & $(XX + YY - 2ZZ)/\sqrt{6}$,\ $(XX - YY)/\sqrt{2}$\\
         \hline
         Symmetric off-diagonal & $3'$ & $(XY + YX)/\sqrt{2}$,\ $(XZ + ZX)/\sqrt{2}$,\ $(YZ + ZY)/\sqrt{2}$\\
         \hline
         Antisymmetric & $3$ & $(XY - YX)/\sqrt{2}$,\ $(XZ - ZX)/\sqrt{2}$,\ $(YZ - ZY)/\sqrt{2}$\\
    \end{tabular}
    \caption{Sub-representations of the $Cl^{\otimes 2}$ group, adapted from
    Eq.~56 of Ref.~\cite{wallman_randomized_2014}. Each row gives the irrep
    label, its dimension, and a spanning set of two-qubit Pauli-basis
    operators.}
    \label{table:irreps}
\end{table}

\begin{table}[t]
    \centering
    \begin{tabular}{c|c|c|c|c|c}
         irrep & $e$ & $(12)$ & $(12)(34)$ & $(123)$ & $(1234)$\\
         \hline
         Identity/Trace & 1 & 1 & 1 & 1 & 1 \\
         \hline
         Left/Right/Antisymmetric & 3 & 0 & $-1$ & 1 & 1 \\
         \hline
         Diagonal traceless & 2 & $-1$ & 2 & 0 & 0 \\
         \hline
         Symmetric off-diagonal & 3 & 0 & $-1$ & $-1$ & $-1$ \\
    \end{tabular}
    \caption{Character table for the irreps of Table~\ref{table:irreps},
    listed against representative conjugacy classes: identity, transposition,
    double transposition, 3-cycle, and 4-cycle.}
    \label{table:character}
\end{table}

\todo{Insert the derivation connecting Tables~\ref{table:irreps} and
\ref{table:character} explicitly (i.e.\ show the character-table computation
step by step, rather than quoting the result). The manuscript SM states this
result without full derivation; a thesis chapter is the appropriate place to
walk through it, likely via Young-tableaux / permutation-representation
arguments.}

Given the block-diagonal structure of Eq.~\eqref{eqn:block_diagonal}, and
further imposing that the twirled channel be trace preserving, a general
twirled two-qubit channel is described by $13$ independent scalars: six
diagonal decay rates,
$\vec\alpha = \{\alpha_L, \alpha_R, \alpha_{tr}, \alpha_{dt}, \alpha_{so},
\alpha_{as}\}$, one per sub-representation, and seven off-diagonal coupling
terms,
$\vec m = \{m_{id\to tr}, m_{L\to R}, m_{R\to L}, m_{L\to as}, m_{R\to as},
m_{as\to L}, m_{as \to R}\}$, describing transitions between multiplicity
spaces of the same irrep. Table~\ref{table:off-diag} associates each
coupling term with a physical error mechanism capable of generating it.

\begin{table}[t]
\centering
\begin{tabular}{c|c|c}
    Coupling & irrep & Mechanism \\
    \hline\hline
    $m_{id\rightarrow tr}$ & 1 & Decays (non-unital) \\
    \hline
    $m_{L\rightarrow R},\ m_{R\rightarrow L}$ & 3 &
    Qubit swap, correlated decays (non-unital) \\
    \hline
    $m_{L\rightarrow as},\ m_{R\rightarrow as},\ m_{as\rightarrow L},\ m_{as\rightarrow R}$
    & 3 & Entanglement \\
\end{tabular}
\caption{Physical mechanisms capable of coupling sub-representation sectors
of the twirled $Cl^{\otimes 2}$ error channel.}
\label{table:off-diag}
\end{table}

This decomposition is a purely group-theoretic statement: it is agnostic to
any particular physical error model and simply enumerates the possible
subspaces that a two-qubit error channel can populate once diagonal
Clifford twirling has been applied. Section~\ref{sec:noise_models} below
specialises this general structure to specific, physically motivated error
channels, and Chapter~\ref{chap:results} uses the resulting scalar decay
rates to build the observable population dynamics measured experimentally.

\todo{Add a subsection discussing the multiplicity structure in more depth
--- i.e.\ what it physically means for an irrep to appear with multiplicity
$>1$, and how this connects to the existence of the coupling terms $\vec m$.
Also worth adding a small worked example reducing a simple noise channel
(e.g.\ pure local depolarisation) through Eq.~\eqref{eqn:block_diagonal} by
hand for pedagogical value.}

% ==========================================================================
\section{Noise models}
\label{sec:noise_models}

Having established the general block-diagonal structure available to any
twirled two-qubit channel, this section reviews the specific classes of
physical error mechanism relevant to the results in
Chapter~\ref{chap:results}, and how each maps onto the sub-representations
of Table~\ref{table:irreps}.

\subsection{Local depolarising noise}

Independent noise acting separately on each qubit is, without loss of
generality (following twirling), well modelled by a depolarising channel,
\begin{equation}
    \rho^{(1q)} \mapsto \mathcal{E}_{dp,a}(\rho^{(1q)})
    = (1-\varepsilon_a)\rho^{(1q)} + \varepsilon_a \frac{I}{2}\mathrm{Tr}(\rho^{(1q)}),
    \label{eqn:1q_depol}
\end{equation}
where the subscript $a \in \{L, R\}$ labels the qubit and $\varepsilon_a$ is
its depolarising error per Clifford. In Liouville form, restricted to the
two-qubit Pauli basis
$\mathcal{P} = \{I\otimes I,\ \vec\sigma \otimes I,\ I \otimes \vec\sigma,\
\vec\sigma \otimes \vec\sigma\}$, independent depolarisation on both qubits
gives
\begin{equation}
    \Lambda_{dp} = \mathrm{diag}^{\mathcal{P}}
    \Bigl(1,\ (1-\varepsilon_L)I_3,\ (1-\varepsilon_R)I_3,\
    (1-\varepsilon_L)(1-\varepsilon_R)I_9\Bigr).
    \label{eqn:local_depol_liouville}
\end{equation}
Local depolarising noise generates no couplings between sub-representations:
it is diagonal in the Schur basis introduced in
Section~\ref{sec:two_copy_cliff}, and drives the two-qubit state toward the
fully mixed state.

\subsection{Correlated (``common'') single-qubit rotation noise}
\label{sec:common_noise}

The central noise model of interest is correlated random-unitary noise, in
which the \emph{same} randomly sampled single-qubit rotation error is
applied to both qubits, here termed \emph{common noise}. This channel is
\begin{equation}
    \rho \mapsto \mathcal{E}^{\mathcal{A}}_{U\otimes U}(\rho)
    = \sum_{a \in \mathcal{A}} p_a\, (U_a \otimes U_a)\, \rho\, (U_a\otimes U_a)^\dagger,
    \label{eqn:rotation_channel}
\end{equation}
where $\mathcal{A}$ is an ensemble of unitaries applied with probabilities
$p_a$ ($\sum_a p_a = 1$), and each $U_a(\theta_a, \hat n_a) =
e^{-i(\theta_a/2)\hat n_a \cdot \hat\sigma}$ is a rotation of angle
$\theta_a$ about axis $\hat n_a$. A subset $\mathcal{A}' \subseteq
\mathcal{A}$ of this class produces genuinely non-separable two-qubit
channels, i.e.\ $\mathcal{E}^{\mathcal{A}'}_{U\otimes U} \neq
\mathcal{E}_U^{(L)} \otimes \mathcal{E}_U^{(R)}$, and it is precisely this
subset that constitutes spatially correlated noise.

Physically, this class captures correlated but non-entangling control and
environmental fluctuations, which are particularly relevant to trapped-ion
and neutral-atom qubits subject to common fluctuations of qubit drive phase,
amplitude, and detuning; magnetic field noise and gradients; beam pointing
and intensity drift; AC Stark-shift fluctuations; motional-state-dependent
over-rotations; and gate-duration calibration errors.

Twirling Eq.~\eqref{eqn:rotation_channel} over $Cl^{\otimes 2}$ reveals that
only the fully symmetric ``trace'' sub-representation is exactly preserved
by common single-qubit rotations, $\alpha_{tr} = 1$, regardless of the
error angle. To second order in a small rotation angle $\theta$, the
remaining scalar decay rates for a single common rotation are
\begin{equation}
    \alpha_L = \alpha_R = \alpha_{as} = 1 - \theta^2/3, \qquad
    \alpha_{tr} = 1, \qquad
    \alpha_{dt} = \alpha_{so} = 1 - \theta^2,
\end{equation}
so that, to leading order, the dynamics depend only on the magnitude of the
error angle and not on its axis. Averaging over the ensemble $\mathcal{A}$
gives $\alpha_i^{(\mathcal{A})} = \sum_{a} p_a \alpha_a$, leading to the
convenient definition of a correlated error per Clifford,
\begin{equation}
    \varepsilon_{\mathrm{corr}} = \frac{1}{3}\sum_{a\in\mathcal{A}} p_a\, \theta_a^2.
    \label{eqn:e_corr}
\end{equation}

\todo{Add the full (non-small-angle) expressions here for completeness, i.e.\
$\alpha_L=\alpha_R=\alpha_{as}=(1+2\cos\theta)/3$ and the corresponding
$\alpha_{dt}$, $\alpha_{so}$ including the anisotropic term $n_{an}$, together
with a short discussion of when the small-angle approximation breaks down.
Also worth adding the tensor-contraction derivation
($s_{tr}, s_{id}, s_{swap}, s_{diag}$) that produces these $\alpha_i$ from
first principles, since it is a reusable technique beyond this specific
channel.}

\subsection{Entangling coherent errors: $ZZ$-type noise}

A distinct and experimentally important source of two-qubit correlation is
coherent entangling coupling, exemplified in superconducting-qubit
architectures by always-on $ZZ$ crosstalk~\cite{zhao_quantum_2022,
fors_comprehensive_2024, alghadeer_crosstalk_2025}, described by the
Hamiltonian
\begin{equation}
    H_{ZZ} = \frac{J_{ZZ}}{2} Z \otimes Z.
\end{equation}
Unlike common single-qubit rotation noise, $H_{ZZ}$ couples the two
multiplicity copies of the $3$-dimensional irrep (the ``left'', ``right'',
and ``antisymmetric'' sub-representations of Table~\ref{table:irreps}), and
so cannot in general be described by a diagonal Liouville matrix in the
Schur basis. A more tractable, stochastic variant of this error,
\begin{equation}
    \rho \mapsto \mathcal{E}_{ZZ}(\rho)
    = (1-p_{ZZ})\rho + p_{ZZ}\, (Z\otimes Z)\rho (Z\otimes Z),
\end{equation}
does remain diagonal, with
\begin{equation}
    \alpha_L = \alpha_R = \alpha_{so} = \alpha_{as} = 1 - 4p_{ZZ}/3,
    \qquad \alpha_{tr} = \alpha_{dt} = 1.
\end{equation}
This channel is included here for completeness and generality; it is not
probed experimentally in this thesis (see Chapter~\ref{chap:results},
Outlook).

\todo{Expand this subsection substantially if $ZZ$-error characterisation
numerics or a proposed experiment are added elsewhere in the thesis. As it
stands this is a very compressed treatment inherited directly from the
manuscript SM; a thesis-length treatment might derive the full coherent
(non-stochastic) block-diagonal Liouville matrix, rather than only the
stochastic special case.}
